\chapter{Face maps, peripheral markings, and the absolute volume}
\label{app:marked-calibration}

This appendix supplies the complete rational dictionary behind
Chapter~\ref{ch:marked}. Its scope is the completed grammar in
Definition~\ref{def:grammar}, with two completed cylindrical root children.
The quadrilateral gluing faces are auxiliary cut data; no physical
four-punctured-boundary representation is inferred from their moduli.
The original closed cut fixes the link group and every peripheral path. At each crossing circle its two shaded triangles meet at their common ideal apex; their union is the configuration called a bowtie. We retain both triangles and their labels. All matrices below are projective: multiplying a matrix by a nonzero scalar does not change its Möbius transformation. All logarithms continue from the strict complete chamber. An unfilled main cusp remains unfilled
after all auxiliary crossing-circle fillings.

\section{Two global incidence checks for the complete anchor}
The complete cells and ordered stack maps are constructed in
Chapter~\ref{ch:geometry}. Two points in that construction deserve an
explicit check before one uses the group presentation.

First, strict output dihedral angles keep the matching quadrilateral a
maximal supported face. Cells on opposite supporting sides intersect only
in that face inside hyperbolic space. Their ideal closures can meet only
on its supporting circle. An additional ideal vertex on that circle would
enlarge the supporting face beyond a quadrilateral edge, contradicting
its strictly nonflat angle. Thus there are no unrecorded ideal coincidences.

Second, at every open stage the four output corners are four unpaired
original terminal singletons. Product consumes $F_2/G_1$ and $F_3/G_0$
and retains $F_0,F_1,G_2,G_3$; turns only permute these singletons.
Consumed pairs never reappear among the exposed vertices of a gluing face. At the root, all four pairs
under $C_j^F\mapsto C_{3-j}^G$ join the disjoint root subtrees. Therefore
every main vertex is exactly a two-terminal pair, not a chain of
identifications, and a leaf's four corners stay distinct. Its two
shaded triangles share only their original apex. An edge and its partner
in the other shaded triangle are distinct, so the complete four-step
cycle contains four distinct edge occurrences. The horospherical links
are positive rectangles, and the canonical shaded-side decoration gives
the actual complete cusp tori described in Chapter~\ref{ch:geometry}.

\section{The selected color action and the complete face-pairing maps}
Write the pure leaf action and its color derivative as
\[
 g(x,t)=J(2x)-2J(x)-J(x-t)-J(x+t),\qquad
 v=g_t=\operatorname{Log}\frac{\sin\pi(x-t)}{\sin\pi(x+t)}.
\]
For physical sign $\eps\in\{0,1,-1\}$, the source-selected shear is
\begin{equation}\label{eq:app-shear}
 \chi=t-\frac{\eps v}{2\pi\ii},\qquad
 f_\eps(x,\chi)=g(x,t)-\frac{\eps v^2}{4\pi\ii},\qquad
 \dd f_\eps=g_x\,\dd x+v\,\dd\chi.
\end{equation}
At completeness its implicit derivative is $1-\eps\ii\cot\pi x\ne0$.
Here \(J(u)=\int_0^u\operatorname{Log}(2\sin\pi s)\,\dd s\), on the
continued complete branch. The finite packet identifies the shear and
its constant. An explicit algebraic check is useful: put
\(E=e^{2\pi\ii x}\), \(H=e^{2\pi\ii\chi}\), and
\(K=e^{-\partial_x f_-}\). The selected source quadratic gives
\begin{equation}\label{eq:app-source-port-curve}
 K^2(1+E)^2+K[(H+H^{-1})E^2+2]+(1-E)^2=0.
\end{equation}
At \(H=1\), it factors as
\((K+1)[K(1+E)^2+(1-E)^2]\), and the source branch is
\(K=\tan^2(\pi x)\). For the negative shear in
\eqref{eq:app-shear}, setting \(T=e^{2\pi\ii t}\) gives
\[
 H=\frac{T(E-T)}{ET-1},\qquad
 T^2+E(H-1)T-H=0,\qquad
 K=\frac{(E-T)(1-ET)}{T(1+E)^2}.
\]
Eliminating \(T\) gives \eqref{eq:app-source-port-curve} with the same
branch. This identifies the \(x\) derivative. The remaining
\(\chi\)-dependent constant is zero by the actual endpoint normalization.
For fixed small \(\chi\), at \(x=1/2\) one has \(v=0\), so the
hidden pure color tends to \(t=\chi\), not to zero. On the continued
branch \(J(1)=J(1/2)=0\) and
\(J(1/2+t)=-J(1/2-t)\), hence \(g(1/2,t)=0\) and the shear
quadratic vanishes. The finite-source action is
\(g(x,\chi)+S(x,\chi)-S(x,0)\), where
\(S(x,\chi)=\int_x^{1/2}\operatorname{Log}\mu(u,\chi)\,\dd u\).
The selected root is regular with \(\mu(1/2,\chi)=1\); both integrals
therefore tend to zero. This verifies the fixed-color constant on the
same holomorphic branch.
Schwarz reflection gives the positive sign. Thus the geometry below
does not select a new source solution or discard a quadratic color phase.

Normalize the ordered gluing face of an elementary pyramid to $(0,\infty,1,R)$, put
$E=e^{2\pi\ii x}$, $T=e^{2\pi\ii t}$, and select
\[
 Z=\frac{1+E}{1+T},\qquad
 R=\frac{T(1+E)^2}{E(1+T)^2},\qquad
 Z^2-(1+E)Z+ER=0.
\]
The two projection charts have apices $Z$ and $Z/E$. Their four maps are listed below. A white map pairs pieces coming from the projection sphere. A completion map is the same algebraic expression at an attachment edge where the adjacent prism supplies the actual face. We distinguish these roles before using the maps in a group relation:
\begin{equation}\label{eq:app-white-matrices}
\begin{aligned}
 W_0&=\begin{pmatrix}1&0\\0&E\end{pmatrix},&
 W_1&=\begin{pmatrix}1&E-1\\0&E\end{pmatrix},\\
 W_2&=\begin{pmatrix}E-R&R(1-E)\\E-1&1-ER\end{pmatrix},&
 W_3&=\begin{pmatrix}R&0\\1-E&ER\end{pmatrix}.
\end{aligned}
\end{equation}
At an elementary pyramid, $W_0,W_2$ are its actual white-face maps. The expressions $W_1,W_3$ become actual maps when the adjacent prism is attached at the edges called cap edges in Chapter~\ref{ch:geometry}. All matrix identities
here are projective. The pure shaded map and the two odd cross-copy maps are
\[
 S=\begin{pmatrix}0&-ER\\1&-1-E\end{pmatrix},\qquad
 A=\begin{pmatrix}R&-R\\1&ER-1-E\end{pmatrix},\qquad
 B=\begin{pmatrix}R/E-1-1/E&R\\-1&R\end{pmatrix}.
\]
They send the original labeled triangle vertices to their prescribed
partners. The apex equation verifies the fourth flag without any
unit-modulus assumption.

For a prism, transport all $W_j$ to the coordinates in
\eqref{eq:ordered-stack}. Direct rational substitution gives
\begin{align}
 W_{F,1}=W_{G,1}=W_{p,1}&=W_{\rm top},&
 W_{F,3}=W_{G,3}=W_{p,3}&=W_{\rm bottom},\notag\\
 W_{F,2}&=W_{G,0},&
 W_{p,0}=W_{F,0},\quad W_{p,2}&=W_{G,2}.
 \label{eq:app-nine-prism-identities}
\end{align}
Here \(p\) denotes the output gluing face, not a filling slope.
The nine identifications are identities of entire maps, not just
cross ratios. They can be verified without geometric continuation:
use independent phase variables
\(\mathsf a=e^{2\pi\ii p}\), \(\mathsf b=e^{2\pi\ii r}\),
\(\mathsf c=e^{2\pi\ii q}\), and the rational coordinates
\[
\begin{aligned}
v_0&=1,\quad v_1=\mathsf a,\quad
v_2=\mathsf a\mathsf b,\quad v_3=\mathsf a\mathsf b\mathsf c,\\
P&=1+\mathsf a-
\frac{\mathsf a[\mathsf a\mathsf b(1+\mathsf c)-1-\mathsf a]}
 {\mathsf a^2\mathsf b^2\mathsf c-\mathsf a}.
\end{aligned}
\]
The three wire phases are
\(\mathsf b\mathsf c,\mathsf a\mathsf b,\mathsf a\mathsf b\mathsf c\).
Normalize each ordered quadrilateral gluing face by
\[
M_{a,b,c}(z)=\frac{(z-a)(c-b)}{(z-b)(c-a)},\qquad
R=M_{a,b,c}(d).
\]
In world coordinates its white map is
\(\overline M^{-1}W_jM\), where the bar replaces
\(\mathsf a,\mathsf b,\mathsf c\) by their reciprocals, as prescribed
by the second chart. Inserting these formulas into
\eqref{eq:app-nine-prism-identities} gives equality entry by entry
up to a common scalar. Thus no unit-circle restriction or missing
translation is hidden in the identities.

The turn
map $T_Q(z)=(R-1)/(z-1)$ and closing map $C_Q(z)=(z-R)/(z-1)$ satisfy
\[
 T_QW_{j+1}(E,R)=W_j(E^{-1},1-R)T_Q,\qquad
 C_QW_j(E^{-1},R)=W_{2-j}(E,R)C_Q.
\]
The identities show that the map on a white face is independent of which constituent cell is used to compute it. Thus every triangular piece of one original white face has the same map, even when its deformed image is not planar. For pure leaves,
writing $\overline S$ for $S$ with $E$ replaced by $E^{-1}$, the edge words are
\[
 W_0S=\overline S W_0,\quad W_2S=\overline S W_2,\quad
 W_3S=\overline S W_1.
\]
For odd leaves they are
\[
 W_0^{-1}BW_2^{-1}A=I,\quad
 W_2^{-1}BW_0^{-1}A=I,\quad
 W_1^{-1}BW_3^{-1}A=I.
\]
All hold for independent nearby $E,R$. They include translation and
dilation, and the reverse words are inverses. The original closed cut
therefore supplies a representation of the actual half-twist partner
group: use its dual two-skeleton presentation, with exactly these white-face seams and shaded four-step relations. No representation is inferred from
the shape parameters of the gluing faces alone. The second charts depend holomorphically on the source
parameters by the reflected-minus convention, rather than by conjugating
those parameters.

A main meridian is $W_{j-1}^{-1}W_j$ at the corresponding corner.
The two maps fix that corner and have equal derivatives there:
\(E^{-1}\) at \(C_0,C_2\), and \(E\) at \(C_1,C_3\), using the
finite coordinate at infinity. Thus its
multiplier is one identically. It is nonidentity at completeness and
stays parabolic nearby. Its commuting peripheral partner is consequently
a translation as well. Every main cusp remains parabolic.

\section{Primitive logs, the absolute regulator, and filling}

We next compute the holonomy of the two actual peripheral curves. Their being primitive is a topological statement about the cusp lattice, so we check their complete translations before using their logarithms. The logarithms below are those of similarity multipliers; they are twice eigenvalue logarithms. This convention explains the factors of two in the volume differential and in the filling equation.
Put
\[
 \rho=\frac{Z}{1+E-Z},\qquad Q=\frac{E-Z}{Z-1}.
\]
The original disk longitude is $W_2^{-1}W_0$, with multiplier $Q^2$ in
the selected flag-at-infinity frame. The pure meridian is $S^{-1}$,
with multiplier $\rho$. The two actual odd meridians are
$A^{-1}W_0$ and $W_0^{-1}B$, with multipliers $\rho Q$ and $\rho/Q$.
Returning to the original complete leaf coordinate, their translations
are $-2$ and $2E$, while the disk longitude is $-2(E+1)$.
The two-rectangle lattice fixes their primitive integral marking and
orientation. The sign is fixed by the original threading convention:
for the negative coupon, the typed loop \([2,1,1,2]\) followed by
\(R^{-1}\) has auxiliary matrix
\[
 M_{\mathrm{aux}}=I-\frac{P+I}{1+E},\qquad \Tr P=-2,
\]
where \(P\) is the actual disk-product matrix and the preferred
projective longitude is \(-P\). Its translation ratio is \(1+E\),
identifying the first path above. Reversing the retained half twist
changes the meridian by one disk longitude and gives the second.
Thus the physical \(R^{-1}\) coupon uses the first path and \(R\)
the second, with no sign inferred solely from a Hessian. Consequently
\[
 \rho=e^{-2\pi\ii t},\qquad Q=e^{-v},\qquad
 U=-2\pi\ii t-\frac{\eps V}{2}=-2\pi\ii\chi,
 \qquad V=-2v.
\]
The half-longitude shear is a relation between the two actual primitive
cusp markings. It does not assert that half a longitude is an integral
loop or replace the actual link by a pure rational filling.
The centered similarity multipliers used for \(U,V\) are squares of
the corresponding eigenvalues. The original lift and pivot character
are retained in the finite source; no additional square-root choice
changes a peripheral word here.

Let $\Phi_\eps(\chi)$ be the selected stationary action in all free wires.
The envelope identity in \eqref{eq:app-shear} gives
\begin{equation}\label{eq:app-primitive}
 U_\ell=-2\pi\ii\chi_\ell,\qquad
 V_\ell=-2\partial_\ell\Phi_\eps,\qquad
 G=2\pi\ii\Phi_\eps,\qquad
 \dd G=\frac12\sum_\ell V_\ell\,\dd U_\ell.
\end{equation}
The source primitive includes the local term $-\eps_\ell V_\ell^2/8$.
It is retained before taking volume. At the complete point,
$\operatorname{Im}G(0)=2\pi F_0$ is the actual half-twist partner volume
from the explicit complete cells and original pairings.

For the nearby framed representation, triangulate the same original cut
balls and retain their peripheral decoration. The ideal-tetrahedron
variation follows directly from the Bloch--Wigner formula
\[
 \dd D(z)=\log|z|\,\dd\arg(1-z)
                    -\log|1-z|\,\dd\arg z.
\]
In the sum over the original oriented tetrahedra, face terms have equal
flags and opposite orientations and cancel. On a peripheral torus,
cut along the recorded primitive \(\mu,\lambda\) to a parallelogram.
The opposite sides cancel except for their two eigenline scale factors;
the remaining alternating product is
\(-2\,\mathcal M\wedge\mathcal L\), where
\(\mathcal M=e^u,\mathcal L=e^v\) and \(u=U/2,v=V/2\).
Here the sign is the boundary orientation of the truncated original
fundamental chain, and the factor two retains both eigenline flags.
Applying the displayed differential to that alternating product gives
\(2(\Re v\,\dd\Im u-\Re u\,\dd\Im v)\).
Different boundary tori have zero mutual intersection, and the
parabolic main tori have zero log contribution. Thus
\[
 \dd\Vol(\rho)
 =\frac12\sum_\ell\bigl(\operatorname{Re}V_\ell\,
       \dd\operatorname{Im}U_\ell
       -\operatorname{Re}U_\ell\,
       \dd\operatorname{Im}V_\ell\bigr).
\]
Using \eqref{eq:app-primitive}, this is the differential of
$\operatorname{Im}G-\frac12\sum\operatorname{Re}U\operatorname{Im}V$.
The complete anchor fixes the constant on the same connected marked
branch. Therefore the absolute identity is
\begin{equation}\label{eq:app-absolute-regulator}
 \Vol(\rho_\chi)=\operatorname{Im}G
       -\frac12\sum_\ell\operatorname{Re}U_\ell\operatorname{Im}V_\ell.
\end{equation}
This uses the original relative fundamental chain and its actual nearby
representation, not equal differentials on unrelated components.

\section{Smooth completion on the actual primitive filled topology}
We give the metric argument used here, including its uniformity in the
independently chosen slopes. Let \(M=Y_E^\eps\) carry the oriented
complete metric constructed above, with holonomy \(\rho_0\). A nearby
\emph{framed} representation consists of the actual face-paired
representation \(\rho\) and, at each cusp, its selected common fixed
flag. In a frame sending that flag to infinity, the two peripheral maps
are complex similarities. Their uniquely small multiplier logs are the
\(U,V\) just computed. We construct one neighborhood of this framed
complete point before choosing the filling slopes.

\subsection{An actual positive metric on the core and its ends}

The first task is to turn nearby holonomy data into local hyperbolic coordinate charts on the same manifold. On a compact truncated part we average small changes of the complete developing map. On each end we use an equivariant affine construction. Both maps are close to the original one on their common compact collar, allowing them to be joined without losing local invertibility.
Choose a compact cusp-truncated core with a slightly larger collar,
finitely many evenly covered balls, and a smooth partition of unity.
Write \(\eta_{j,\gamma}\) for its lifts, with one chosen lift of each
ball and \(\gamma\in\pi_1(M)\). In the future unit hyperboloid model
of hyperbolic space put
\[
 Y_\rho(x)=\sum_{j,\gamma}\eta_{j,\gamma}(x)
                 \rho(\gamma)\rho_0(\gamma)^{-1}x,
 \qquad
 E_\rho(x)=\frac{Y_\rho(x)}
                   {\sqrt{-\langle Y_\rho(x),Y_\rho(x)\rangle}}.
\]
Every summand is future timelike; their convex average is future
timelike too, so the normalization is defined. Reindexing the sum proves
\(E_\rho(\rho_0(\delta)x)=\rho(\delta)E_\rho(x)\).
Only finitely many terms occur on a compact fundamental part, and at
\(\rho_0\) the map is the identity. Thus it is \(C^1\)-close to the
identity and is an orientation-preserving local diffeomorphism in one
fixed representation neighborhood. This does not assume that the
deformed action is already discrete or free.

For one complete cusp write its universal coordinates as
\((w,s)\in\CC\times[s_0,\infty)\), with translation lattice
\(\Lambda\), and complete development \((w,e^s)\) in upper
half-space. Denote the deformed affine peripheral action by
\(A_\lambda(z)=a_\lambda z+b_\lambda\). Choose a nonnegative compactly
supported smooth \(\eta\) with
\(\sum_{\lambda\in\Lambda}\eta(w-\lambda)=1\), and define
\[
 d_\rho(w)=\sum_{\lambda\in\Lambda}\eta(w-\lambda)
                              A_\lambda(w-\lambda).
\]
Then \(d_\rho(w+\lambda)=A_\lambda d_\rho(w)\). On a fixed closed
fundamental parallelogram it is \(C^1\)-close to \(w\), so its real
Jacobian is positive; equivariance propagates that fact everywhere.
Let \(\xi_\rho\) be the real-linear extension of
\(\lambda\mapsto\log|a_\lambda|\). The map
\[
 F_\rho(w,s)=\bigl(d_\rho(w),e^{s+\xi_\rho(w)}\bigr)
\]
is exactly peripheral-equivariant. Its differential has positive
determinant \(e^{s+\xi_\rho(w)}\det_\RR Dd_\rho(w)\) for every
\(s\ge s_0\). Restore the moving flag frame and extend to all cusp
lifts by equivariance. Geodesically interpolate \(E_\rho\) and
\(F_\rho\) on the compact collar with an invariant smooth cutoff.
There both maps are \(C^1\)-close to the complete development, so the
interpolation remains a positive local diffeomorphism after one further
uniform shrink. Pullback of the hyperbolic metric now gives an actual
oriented nonsingular metric on the original manifold with its ideal
ends deleted.

\subsection{The logarithmic link and its primitive angular period}

At a filled cusp, the nearby peripheral similarities have two common fixed points. Sending one to infinity and translating the other to zero makes them multiplications. Their logarithms then act by translations on a Euclidean plane. We must show that the translations are exactly the selected small logs $U,V$, with no unnoticed multiple of $2\pi\ii$. Only after this check may the large integer $p$ be used in the relation $U+pV=-2\pi\ii$.
At a cusp to be filled, the equation \(U+pV=-2\pi\ii\) ensures that
at least one multiplier differs from one. Commuting affine maps then
have a common finite fixed point \(q\), in addition to infinity. For
either generator with \(a\ne1\), \(q=b/(1-a)\). Its numerator stays
away from zero near completeness, whereas both generator multipliers
tend to one; hence \(|q|\to\infty\) uniformly, even if one multiplier
equals one. The developed fundamental parallelogram and fixed generator
paths have uniformly bounded images and lengths. After shrinking the
neighborhood they miss \(q\), and equivariance implies that the entire
development misses \(q\).

On the simply connected link cover choose
\(f=\operatorname{Log}(d_\rho-q)\). Its differential is invertible
as a real map and is invariant under deck transformations. Consequently
\(|\dd f|^2\) descends to a nonsingular Euclidean metric on the compact
torus, hence is complete. Lifting Euclidean segments proves that its
developing local isometry \(f\) covers \(\CC\); because the plane is
simply connected, this is a global isometry. Its periods form an
injective real rank-two lattice. A fixed generator path of developed
length at most \(L\), contained in \(|d_\rho|\le R\), satisfies
\[
 |\Delta f|\le\frac{L}{|q|-R}\longrightarrow0.
\]
Exponentiating its increment gives the corresponding similarity
multiplier. The increment and its prescribed small log therefore differ
by an integer multiple of \(2\pi\ii\) which is exactly zero near
completeness. Thus the actual two log-link periods are \(U,V\),
before they are multiplied by any large slope. This is why the same
representation neighborhood works for all slopes.

Translate \(q\) to zero. On the compact link quotient the function
\(e^{s_0+\xi_\rho(w)}/|d_\rho(w)-q|\) is bounded above. The explicit
end therefore contains an entire equivariant epigraph
\(y>A|d_\rho-q|\). With
\[
 r=\frac{y}{|d_\rho-q|},\qquad
 R=\operatorname{arsinh}(1/r),\qquad
 \theta=\Im f,\qquad
 X=\Re f+\tfrac12\log(1+r^2),
\]
direct substitution into the upper-half-space metric gives
\begin{equation}\label{eq:app-smooth-tube}
 \dd s^2=\dd R^2+\cosh^2R\,\dd X^2+
                           \sinh^2R\,\dd\theta^2.
\end{equation}
The epigraph contains a whole punctured tube \(0<R<R_0\). A log period
\(L\) acts by \((X,\theta)\mapsto(X+\Re L,\theta+\Im L)\).
The actual primitive curve \(\mu+p\lambda\) has period
\(-2\pi\ii\), exactly one angular turn. Its complementary primitive
period can be taken to be \(V\), which has nonzero real part since
\(U,V\) are real-linearly independent. Adding the axis \(R=0\) in
\eqref{eq:app-smooth-tube} is therefore a smooth hyperbolic solid torus,
whose meridian is exactly \(\mu+p\lambda\). Its affine holonomy is
the identity, since it fixes both \(q\) and infinity with multiplier
one. The representation descends to the actual filled group. Reversing
the angular sign changes the parametrization of the meridian, not the
orientation of the metric.

At an unfilled main cusp both multipliers are one. The same construction
gives a translation-equivariant Euclidean link development, globally
the plane by completeness of the compact link; \(\xi_\rho=0\).
Its end contains a full horoball quotient, so it is complete and has
finite volume. The remaining region is compact. Attaching the finitely
many bounded-radius solid tori therefore gives a complete finite-volume
hyperbolic metric on precisely the filled topology. Its complete
universal developing map covers simply connected hyperbolic space,
so the holonomy is discrete and faithful. This proof used one fixed
framed neighborhood and no comparison between different slope sizes.

\subsection{The source saddle and its absolute knot-complement volume}
The full source filling transform selects a critical point of
\[
 S_{\eps,\mathbf p}(\chi)=\Phi_\eps(\chi)
       +\frac{\pi\ii}{2}\sum_\ell\frac{(\chi_\ell-1)^2}{p_\ell}.
\]
Its equations are exactly $U_\ell+p_\ell V_\ell=-2\pi\ii$.
For existence set \(r_\ell=1/p_\ell\) and write them as
\(V_\ell+r_\ell U_\ell=-2\pi\ii r_\ell\). At \(r=0\) their
\(\chi\)-Jacobian is \(-4\pi\ii\,\diag(\tau_\ell)\), with
\(\Im\tau_\ell>0\), by the source packet. It is invertible.
The holomorphic implicit-function theorem thus gives one nearby solution
for all \(r\) in a fixed polydisk, independently of their relative
sizes and signs. It lies in the framed neighborhood just constructed
after that polydisk is shrunk once more.
For independently sufficiently long fixed signed integral $p_\ell$, the
marked geometric completion gives smooth angle $2\pi$ on these primitive
filled curves, with the unfilled main cusp parabolic. The physical leaf
power is $R^{\eps_\ell-2p_\ell}$, once, in its original disk framing.
At this saddle, \eqref{eq:app-absolute-regulator} is
\[
 \Vol\bigl(S^3\setminus K_{F,G,\eps}(\mathbf p)\bigr)
  =2\pi\operatorname{Re}S_{\eps,\mathbf p}(\chi_*).
\]
All auxiliary circles have been filled and the chosen main component is
the knot. The knot cusp is unfilled. The separate full-source alias,
complex-packet and entire Fourier-period arguments are what identify this
geometric action with the original invariant's limiting rate.
